% !PRINTED 253
\weberpage{253}
%<<CONT
$j(\omega)$ nach Potenzen von $q$ den Koeffizienten 1 hat, schließt man
leicht, daß die numerischen Koeffizienten in diesen Gleichungen
\gesperrt{rationale ganze Zahlen} sind.

Für $p=5$ hat die fragliche Gleichung die Form:
\[
x^6 + A_1\gamma_2^2x^5 + A_2\gamma_2x^4 + A_3x^3 + A_4\gamma_2x + 5 = 0,
\]
worin die $A_1$, $A_2$, $A_3$, $A_4$ ganze rationale Funktionen von
$j(\omega)$ und, wie leicht zu sehen, Konstante sind.

Aus den Anfangsgliedern ergibt sich sofort
\[
A_1 = 0, \quad A_2 = 0, \quad A_4 = -1.
\]

Um aber $A_3$ zu bestimmen, geht man in der Entwickelung von $x_0$ bis
zum nächstfolgenden Gliede:
\[
x_0 = q^{-\frac{7}{15}}\Bigl(1 - 2q^{\frac{2}{5}}\dots\Bigr),
\]
woraus man $A_3 = 10$ erhält; also lautet für $n=5$ die
Multiplikatorgleichung
\weq{8}{x^6 + 10x^3 - \gamma_2x + 5 = 0.}

In gleicher Weise berechnet man die Gleichungen für $p=7$, $p=11$,
indem man die Entwickelungen von $x_0$ benutzt:
\[
p=7:\ x_0 = -q^{-\frac{1}{7}}\Bigl(1 - 2q^{\frac{2}{7}} - q^{\frac{4}{7}} + 2q^{\frac{6}{7}}\dots\Bigr),
\]
\[
p=11:\ x_0 = -q^{-\frac{5}{33}}\Bigl(1 - 2q^{\frac{2}{11}} - q^{\frac{4}{11}} + 2q^{\frac{6}{11}} + q^{\frac{8}{11}} + 2q^{\frac{10}{11}} + \dots\Bigr),
\]
während von $\gamma_2(\omega)$, $\gamma_3(\omega)$ immer nur die ersten
Glieder gebraucht werden. So findet sich
\weq{9}{p=7:\quad x^8 + 7\cdot2x^6 + 7\cdot9x^4 + 7\cdot10x^2 + \gamma_3x - 7 = 0,}
\weq{10}{p=11:\quad x^{12} - 11\cdot90x^6 + 11\cdot40\gamma_2x^4 + 11\cdot15\gamma_3x^3 + 11\cdot2\gamma_2^2x^2 + \gamma_2\gamma_3x - 11 = 0.}

Wenn $n$ eine ungerade Quadratzahl ist, so nehmen wir $c$ \gesperrt{durch 8
teilbar} an. Es sind dann $a:e$ und $\partial:e$ ebenfalls
Quadratzahlen. Es ist ferner
\weq{11}{\alpha \equiv \delta \equiv 0\ (\mathrm{mod}\ 8), \quad \lambda \equiv 1\ (\mathrm{mod}\ 8)}
\weq{12}{\left.
\begin{aligned}
a &= \beta\partial_2, & a_2 &= \beta\partial, & a &\equiv \partial \equiv e \equiv 1\\
c &= \delta\partial_2, & c_2 &= -\alpha\partial, & a_2 &\equiv \partial_2 \equiv e_2 \equiv 1
\end{aligned}
\right\}\ (\mathrm{mod}\ 8).}

Die Einheitswurzel $r$ in (2) erhält den Wert
\weq{13}{r = \left(\frac{\alpha}{\beta}\right) i^{\frac{1-\beta}{2}} e^{\frac{2\pi i}{3}\frac{\beta(\alpha+\delta) - (\beta^2-1)\alpha\gamma}{8}}.}
